\documentclass[aspectratio=169,11pt]{beamer}

% Tema e cores
\usetheme{Madrid}
\usecolortheme{whale}
\setbeamertemplate{navigation symbols}{}
\setbeamertemplate{footline}[frame number]

% Pacotes
\usepackage[utf8]{inputenc}
\usepackage[T1]{fontenc}
\usepackage[portuguese]{babel}
\usepackage{graphicx}
\usepackage{booktabs}
\usepackage{amsmath}
\usepackage{xcolor}
\usepackage{tikz}
\usepackage{multirow}

% Cores customizadas
\definecolor{lnccblue}{RGB}{52, 152, 219}
\definecolor{alertred}{RGB}{231, 76, 60}
\definecolor{successgreen}{RGB}{39, 174, 96}
\definecolor{warningorange}{RGB}{243, 156, 18}

% Configurações
\setbeamercolor{title}{fg=white,bg=lnccblue}
\setbeamercolor{frametitle}{fg=white,bg=lnccblue}
\setbeamercolor{block title}{bg=lnccblue,fg=white}
\setbeamercolor{block body}{bg=lnccblue!10}
\setbeamercolor{block title alerted}{bg=alertred,fg=white}
\setbeamercolor{block body alerted}{bg=alertred!10}

% Informações do documento
\title[Análise de Vazamento de Dados]{Por que Representações Simples Apresentam\\Performance Elevada?}
\subtitle{Análise de Vazamento de Dados em Predição de\\Atividade Quinase-Composto}
\author{Análise Exploratória de Dados}
\institute{LNCC - Laboratório Nacional de Computação Científica}
\date{\today}

\begin{document}

% =============================================================================
% SLIDE DE TÍTULO
% =============================================================================
\begin{frame}
    \titlepage
\end{frame}

% =============================================================================
% SUMÁRIO
% =============================================================================
\begin{frame}{Sumário}
    \tableofcontents
\end{frame}

% =============================================================================
% SEÇÃO 1: CONTEXTO E PROBLEMA
% =============================================================================
\section{Contexto e Problema}

\begin{frame}{Contexto: Predição de Atividade Quinase-Composto}
    \begin{columns}
        \begin{column}{0.5\textwidth}
            \textbf{Objetivo:}
            \begin{itemize}
                \item Predizer se um composto químico é \textbf{ativo} ou \textbf{inativo} contra uma quinase
                \item Classificação binária: ativo ($\leq 1\mu M$) vs inativo
            \end{itemize}

            \vspace{0.5cm}
            \textbf{Representações Utilizadas:}
            \begin{itemize}
                \item \textbf{Quinases:} One-hot encoding
                \item \textbf{Compostos:} Morgan Fingerprints (2048 bits)
            \end{itemize}
        \end{column}
        \begin{column}{0.5\textwidth}
            \begin{block}{Características do Dataset}
                \begin{itemize}
                    \item 15.616 pares composto-quinase
                    \item 8.131 compostos únicos
                    \item 231 quinases únicas
                    \item 56,8\% ativos / 43,2\% inativos
                \end{itemize}
            \end{block}

            \vspace{0.3cm}
            \begin{block}{Split Utilizado}
                80\% treino / 10\% validação / 10\% teste\\
                (Split aleatório estratificado)
            \end{block}
        \end{column}
    \end{columns}
\end{frame}

\begin{frame}{O Problema: Performance Inesperadamente Alta}
    \begin{alertblock}{Resultados Observados}
        \centering
        \begin{tabular}{lcc}
            \toprule
            \textbf{Modelo} & \textbf{Accuracy} & \textbf{MCC} \\
            \midrule
            KNN & 0.876 & 0.747 \\
            MLP & 0.887 & 0.768 \\
            \bottomrule
        \end{tabular}
    \end{alertblock}

    \vspace{0.5cm}
    \begin{center}
        \Large\textcolor{alertred}{\textbf{Questão Central:}}\\[0.3cm]
        \textit{Por que uma representação tão simples consegue\\performance tão elevada?}\\[0.5cm]
        \textcolor{warningorange}{\textbf{Expectativa: O problema deveria ser muito mais difícil!}}
    \end{center}
\end{frame}

% =============================================================================
% SEÇÃO 2: VAZAMENTO DE DADOS
% =============================================================================
\section{Análise de Vazamento de Dados}

\begin{frame}{Problema Crítico: Vazamento de Compostos}
    \begin{columns}
        \begin{column}{0.55\textwidth}
            \includegraphics[width=\textwidth]{01_leakage_analysis.png}
        \end{column}
        \begin{column}{0.45\textwidth}
            \begin{alertblock}{Descoberta Alarmante}
                \begin{itemize}
                    \item \textbf{64,7\%} das linhas de teste têm composto vazado
                    \item \textbf{36,6\%} são duplicatas exatas (mesmo composto + quinase)
                \end{itemize}
            \end{alertblock}

            \vspace{0.2cm}
            \begin{block}{Detalhamento}
                \small
                \begin{itemize}
                    \item 820 de 1.355 compostos únicos do teste (60,5\%) aparecem no treino
                    \item Esses 820 compostos geram 64,7\% das linhas
                \end{itemize}
            \end{block}

            \vspace{0.2cm}
            \begin{block}{Consequência}
                O modelo pode estar \textbf{lembrando} ao invés de \textbf{generalizando}!
            \end{block}
        \end{column}
    \end{columns}
\end{frame}

% =============================================================================
% SEÇÃO 3: BASELINE LOOKUP
% =============================================================================
\section{Modelo Lookup Baseline}

\begin{frame}{Hipótese: Modelo Simples de ``Lookup''}
    \textbf{Ideia:} Se os modelos estão memorizando, um modelo que apenas ``olha'' a classe majoritária deveria ter performance similar.

    \vspace{0.5cm}
    \begin{block}{Modelos Lookup Implementados}
        \begin{enumerate}
            \item \textbf{Lookup por Composto:} Prediz a classe majoritária do composto no treino
            \item \textbf{Lookup por Quinase:} Prediz a classe majoritária da quinase no treino
            \item \textbf{Lookup Composto + Quinase:} Combina ambas as estratégias
        \end{enumerate}
    \end{block}

    \vspace{0.3cm}
    \begin{center}
        \textit{Se o composto/quinase não existe no treino, prediz a classe global majoritária.}
    \end{center}
\end{frame}

\begin{frame}{Resultados: Lookup vs KNN}
    \begin{center}
        \includegraphics[width=0.85\textwidth]{02_baseline_comparison.png}
    \end{center}
\end{frame}

\begin{frame}{Análise dos Resultados Lookup}
    \begin{columns}
        \begin{column}{0.5\textwidth}
            \begin{table}
                \centering
                \small
                \begin{tabular}{lcc}
                    \toprule
                    \textbf{Modelo} & \textbf{Acc} & \textbf{MCC} \\
                    \midrule
                    Lookup: Composto & 0.787 & 0.577 \\
                    Lookup: Quinase & 0.750 & 0.510 \\
                    Lookup: Comp+Quin & \textbf{0.862} & \textbf{0.717} \\
                    \midrule
                    KNN (Original) & 0.876 & 0.747 \\
                    \bottomrule
                \end{tabular}
            \end{table}
        \end{column}
        \begin{column}{0.5\textwidth}
            \begin{alertblock}{Conclusão}
                O modelo \textbf{Lookup simples} (apenas olha classe majoritária) tem performance \textbf{comparável} ao KNN!
            \end{alertblock}

            \vspace{0.3cm}
            \begin{block}{Implicação}
                \textcolor{alertred}{\textbf{O KNN não está aprendendo padrões complexos --- está memorizando!}}
            \end{block}
        \end{column}
    \end{columns}
\end{frame}

% =============================================================================
% SEÇÃO 4: DESBALANCEAMENTO POR QUINASE
% =============================================================================
\section{Desbalanceamento por Quinase}

\begin{frame}{Distribuição de Classes por Quinase}
    \begin{center}
        \includegraphics[width=0.9\textwidth]{03_kinase_imbalance.png}
    \end{center}
\end{frame}

\begin{frame}{Quinases Trivialmente Previsíveis}
    \begin{columns}
        \begin{column}{0.5\textwidth}
            \begin{alertblock}{Estatísticas Alarmantes}
                \begin{itemize}
                    \item \textbf{70,6\%} das quinases são desbalanceadas ($>$80\% uma classe)
                    \item \textbf{45,6\%} das linhas vêm de quinases extremamente desbalanceadas
                \end{itemize}
            \end{alertblock}

            \vspace{0.3cm}
            \begin{block}{Exemplos Extremos}
                \small
                \begin{itemize}
                    \item Glycogen synthase kinase-3 beta: \textbf{100\% ativos}
                    \item Muitas quinases: \textbf{100\% inativos}
                \end{itemize}
            \end{block}
        \end{column}
        \begin{column}{0.5\textwidth}
            \begin{block}{Consequência}
                Para muitas quinases, o modelo pode simplesmente predizer a classe majoritária e ter alta acurácia!
            \end{block}

            \vspace{0.5cm}
            \textbf{Exemplo:}\\
            Se uma quinase tem 100\% de compostos inativos, o modelo precisa apenas ``aprender'' que essa quinase $\rightarrow$ inativo.
        \end{column}
    \end{columns}
\end{frame}

% =============================================================================
% SEÇÃO 5: CONSISTÊNCIA DOS COMPOSTOS
% =============================================================================
\section{Consistência de Classe dos Compostos}

\begin{frame}{Comportamento dos Compostos através de Quinases}
    \begin{center}
        \includegraphics[width=0.9\textwidth]{04_compound_consistency.png}
    \end{center}
\end{frame}

\begin{frame}{Análise de Consistência}
    \begin{columns}
        \begin{column}{0.5\textwidth}
            \begin{table}
                \centering
                \begin{tabular}{lc}
                    \toprule
                    \textbf{Categoria} & \textbf{\%} \\
                    \midrule
                    Uma quinase apenas & 80,2\% \\
                    Perfeitamente consistente & 15,2\% \\
                    Inconsistente & 4,6\% \\
                    \bottomrule
                \end{tabular}
            \end{table}

            \vspace{0.3cm}
            \begin{block}{Multi-quinase}
                \textbf{76,7\%} dos compostos multi-quinase são perfeitamente consistentes!
            \end{block}
        \end{column}
        \begin{column}{0.5\textwidth}
            \begin{alertblock}{Implicação Crítica}
                O \textbf{fingerprint do composto sozinho} carrega a maior parte do sinal preditivo!
            \end{alertblock}

            \vspace{0.3cm}
            \begin{center}
                \textit{Ou seja, só de olhar para o composto, já é possível saber se ele é ativo ou inativo para a maioria das quinases.}
            \end{center}
        \end{column}
    \end{columns}
\end{frame}

% =============================================================================
% SEÇÃO 6: SIMILARIDADE QUÍMICA
% =============================================================================
\section{Similaridade Química}

\begin{frame}{Similaridade Tanimoto: Teste vs Treino}
    \begin{center}
        \includegraphics[width=0.9\textwidth]{05_similarity_analysis.png}
    \end{center}
\end{frame}

\begin{frame}{Por que o KNN Funciona?}
    \begin{columns}
        \begin{column}{0.5\textwidth}
            \begin{table}
                \centering
                \begin{tabular}{lc}
                    \toprule
                    \textbf{Similaridade} & \textbf{\%} \\
                    \midrule
                    Muito similar ($>$0.8) & 43,0\% \\
                    Similar (0.6--0.8) & 46,8\% \\
                    Moderada (0.4--0.6) & 5,8\% \\
                    Baixa ($<$0.4) & 4,4\% \\
                    \bottomrule
                \end{tabular}
            \end{table}
        \end{column}
        \begin{column}{0.5\textwidth}
            \begin{alertblock}{Descoberta}
                \textbf{$\sim$90\%} dos compostos ``novos'' no teste têm vizinho muito próximo no treino!
            \end{alertblock}

            \vspace{0.3cm}
            \begin{block}{Conclusão}
                O KNN funciona bem porque a maioria dos compostos de teste é \textbf{quimicamente muito similar} aos compostos de treino.
            \end{block}
        \end{column}
    \end{columns}
\end{frame}

% =============================================================================
% SEÇÃO 7: COMPARAÇÃO DE SPLITS
% =============================================================================
\section{Comparação de Cenários de Avaliação}

\begin{frame}{Três Cenários de Split}
    \begin{block}{Cenário 1: Split Aleatório (Original)}
        Split 80/10/10 estratificado --- \textcolor{alertred}{COM vazamento de dados}
    \end{block}

    \begin{block}{Cenário 2: Split por Composto}
        Nenhum composto do teste aparece no treino --- \textcolor{warningorange}{Sem vazamento de compostos}
    \end{block}

    \begin{block}{Cenário 3: Novo Composto + Nova Quinase}
        Compostos E quinases do teste nunca foram vistos no treino --- \textcolor{successgreen}{\textbf{Generalização Verdadeira}}
    \end{block}

    \vspace{0.3cm}
    \begin{center}
        \textit{O Cenário 3 representa o caso de uso real: predizer atividade para quinases e compostos completamente novos.}
    \end{center}
\end{frame}

\begin{frame}{Resultados por Cenário de Avaliação}
    \begin{center}
        \includegraphics[width=0.95\textwidth]{06_split_comparison.png}
    \end{center}
\end{frame}

\begin{frame}{Performance Inflada vs Performance Real}
    \begin{center}
        \includegraphics[width=0.95\textwidth]{07_inflated_vs_real_performance.png}
    \end{center}
\end{frame}

\begin{frame}{Queda Dramática de Performance}
    \begin{columns}
        \begin{column}{0.5\textwidth}
            \begin{table}
                \centering
                \begin{tabular}{lccc}
                    \toprule
                    & \textbf{Reportado} & \textbf{Real} & \textbf{Queda} \\
                    \midrule
                    \multicolumn{4}{c}{\textbf{MCC}} \\
                    \midrule
                    KNN & 0.747 & 0.407 & \textcolor{alertred}{-45\%} \\
                    MLP & 0.768 & 0.275 & \textcolor{alertred}{-64\%} \\
                    \midrule
                    \multicolumn{4}{c}{\textbf{Accuracy}} \\
                    \midrule
                    KNN & 0.876 & 0.680 & \textcolor{alertred}{-22\%} \\
                    MLP & 0.887 & 0.607 & \textcolor{alertred}{-32\%} \\
                    \bottomrule
                \end{tabular}
            \end{table}
        \end{column}
        \begin{column}{0.5\textwidth}
            \begin{alertblock}{Conclusão}
                A queda \textbf{dramática} de performance confirma que os modelos estavam \textbf{MEMORIZANDO}, não aprendendo padrões generalizáveis!
            \end{alertblock}

            \vspace{0.3cm}
            \begin{block}{Nota}
                O MLP sofre mais no cenário de generalização verdadeira (MCC cai para 0.275 --- quase aleatório!)
            \end{block}
        \end{column}
    \end{columns}
\end{frame}

% =============================================================================
% SEÇÃO 8: CONCLUSÕES
% =============================================================================
\section{Conclusões e Implicações}

\begin{frame}{Resumo das Descobertas}
    \begin{enumerate}
        \item \textbf{Vazamento Massivo:} 60,5\% dos compostos de teste aparecem no treino
        \vspace{0.3cm}
        \item \textbf{Desbalanceamento Extremo:} 70,6\% das quinases são trivialmente previsíveis
        \vspace{0.3cm}
        \item \textbf{Consistência de Compostos:} O fingerprint sozinho carrega o sinal preditivo
        \vspace{0.3cm}
        \item \textbf{Alta Similaridade:} 90\% dos compostos de teste têm vizinho muito próximo
        \vspace{0.3cm}
        \item \textbf{Performance Inflada:} MCC cai de 0.77 para 0.27 em generalização real
    \end{enumerate}
\end{frame}

\begin{frame}{Implicações para a Área}
    \begin{alertblock}{Crítica Central}
        \textbf{Modelos de triagem de quinases que não tiveram curadoria adequada dos dados podem estar superestimando drasticamente sua capacidade de generalização!}
    \end{alertblock}

    \vspace{0.5cm}
    \begin{block}{Recomendações}
        \begin{itemize}
            \item Sempre verificar vazamento de compostos entre splits
            \item Avaliar desbalanceamento por quinase
            \item Usar splits que garantam generalização verdadeira
            \item Reportar métricas em cenários de compostos E quinases novos
            \item Comparar com baselines simples (Lookup)
        \end{itemize}
    \end{block}
\end{frame}

\begin{frame}{Conclusão Final}
    \begin{center}
        \Large
        \textcolor{lnccblue}{\textbf{Os modelos (KNN, MLP) NÃO estão aprendendo padrões complexos de interação quinase-composto.}}

        \vspace{1cm}

        \textcolor{alertred}{\textbf{Eles estão MEMORIZANDO devido a:}}

        \vspace{0.5cm}
        \begin{enumerate}
            \item Vazamento massivo de compostos
            \item Desbalanceamento extremo por quinase
            \item Comportamento consistente dos compostos
            \item Alta similaridade química ao treino
        \end{enumerate}

        \vspace{1cm}
        \textit{Esta análise é fundamental para validar qualquer modelo de predição de atividade quinase-composto.}
    \end{center}
\end{frame}

% =============================================================================
% SLIDE FINAL
% =============================================================================
\begin{frame}
    \begin{center}
        \Huge\textcolor{lnccblue}{\textbf{Obrigado!}}

        \vspace{1cm}
        \Large
        Perguntas?

        \vspace{1.5cm}
        \normalsize
        \textit{LNCC - Laboratório Nacional de Computação Científica}\\
        \textit{Análise Exploratória de Dados}
    \end{center}
\end{frame}

\end{document}
